\documentclass[a4paper,10pt,BCOR10mm,oneside,headsepline]{scrartcl}
\usepackage[portuges]{babel}
\usepackage[utf8]{inputenc}
\usepackage{amsmath}
\usepackage{wasysym}% provides \ocircle and \Box
\usepackage{enumitem}% easy control of topsep and leftmargin for lists
\usepackage{color}% used for background color
\usepackage{forloop}% used for \Qrating and \Qlines
\usepackage{ifthen}% used for \Qitem and \QItem
\usepackage{typearea}
\usepackage{graphicx}
\usepackage{booktabs}
\usepackage{listings}
\usepackage{color}
\usepackage{algorithm}
\usepackage{algorithmic}
\usepackage{amssymb}
\usepackage{minted}
\usepackage{mdframed}
\usepackage{wasysym}
\usepackage{xcolor}

\usepackage{tgtermes}

\areaset{17cm}{26cm}
\setlength{\topmargin}{-2cm}

\definecolor{codegreen}{rgb}{0,0.6,0}
\definecolor{codegray}{rgb}{0.5,0.5,0.5}
\definecolor{codepurple}{rgb}{0.58,0,0.82}
\definecolor{backcolour}{rgb}{0.95,0.95,0.92}
 
\lstdefinestyle{mystyle}{
    backgroundcolor=\color{backcolour},   
    commentstyle=\color{codegreen},
    keywordstyle=\color{magenta},
    numberstyle=\tiny\color{codegray},
    stringstyle=\color{codepurple},
    basicstyle=\footnotesize,
    breakatwhitespace=false,         
    breaklines=true,                 
    captionpos=b,                    
    keepspaces=true,                 
    numbers=left,                    
    numbersep=5pt,                  
    showspaces=false,                
    showstringspaces=false,
    showtabs=false,                  
    tabsize=2
}
 
\lstset{style=mystyle}

\mdfsetup{
    skipabove=12pt,
    skipbelow=12pt,
    innertopmargin=8pt,
    innerbottommargin=8pt,
    innerleftmargin=10pt,
    innerrightmargin=10pt,
    splittopskip=12pt,
    splitbottomskip=12pt,
}


\algsetup{linenodelimiter=.,linenosize=\tiny}

\renewcommand{\algorithmicrequire}{\textbf{Entrada:}}
\renewcommand{\algorithmicensure}{\textbf{Saída:}}
\floatname{algorithm}{Algoritmo}

%% \Qq = Questionaire question. Oh, this is just too simple. It helps
%% making it easy to globally change the appearance of questions.
\newcommand{\Qq}[1]{\textbf{#1}}
%\newcommand{\Qq}[1]{{#1}}

%% \QO = Circle or box to be ticked. Used both by direct call and by
%% \Qrating and \Qlist.
\newcommand{\QO}{$\Box$}% or: $\ocircle$

%% \Qrating = Automatically create a rating scale with NUM steps, like
%% this: 0--0--0--0--0.
\newcounter{qr}
\newcommand{\Qrating}[1]{\QO\forloop{qr}{1}{\value{qr} < #1}{---\QO}}

%% \Qline = Again, this is very simple. It helps setting the line
%% thickness globally. Used both by direct call and by \Qlines.
\newcommand{\Qline}[1]{\noindent\rule{#1}{0.6pt}}

%% \Qlines = Insert NUM lines with width=\linewith. You can change the
%% \vskip value to adjust the spacing.
\newcounter{ql}
\newcommand{\Qlines}[1]{\forloop{ql}{0}{\value{ql}<#1}{\vskip0em\Qline{\linewidth}}}

%% \Qlist = This is an environment very similar to itemize but with
%% \QO in front of each list item. Useful for classical multiple
%% choice. Change leftmargin and topsep accourding to your taste.
\newenvironment{Qlist}{%
\renewcommand{\labelitemi}{\QO}
\begin{itemize}[leftmargin=1.5em,topsep=-.5em]
}{%
\end{itemize}
}

%% \Qtab = A "tabulator simulation". The first argument is the
%% distance from the left margin. The second argument is content which
%% is indented within the current row.
\newlength{\qt}
\newcommand{\Qtab}[2]{
\setlength{\qt}{\linewidth}
\addtolength{\qt}{-#1}
\hfill\parbox[t]{\qt}{\raggedright #2}
}

%% \Qitem = Item with automatic numbering. The first optional argument
%% can be used to create sub-items like 2a, 2b, 2c, ... The item
%% number is increased if the first argument is omitted or equals 'a'.
%% You will have to adjust this if you prefer a different numbering
%% scheme. Adjust topsep and leftmargin as needed.
\newcounter{itemnummer}
\newcommand{\Qitem}[2][]{% #1 optional, #2 notwendig
\ifthenelse{\equal{#1}{}}{\stepcounter{itemnummer}}{}
\ifthenelse{\equal{#1}{a}}{\stepcounter{itemnummer}}{}
\begin{enumerate}[topsep=2pt,leftmargin=2.8em]
\item[\textbf{\arabic{itemnummer}#1.}] #2
\end{enumerate}
}

%% \QItem = Like \Qitem but with alternating background color. This
%% might be error prone as I hard-coded some lengths (-5.25pt and
%% -3pt)! I do not yet understand why I need them.
\definecolor{bgodd}{rgb}{0.8,0.8,0.8}
\definecolor{bgeven}{rgb}{0.9,0.9,0.9}
\newcounter{itemoddeven}
\newlength{\gb}
\newcommand{\QItem}[2][]{% #1 optional, #2 notwendig
\setlength{\gb}{\linewidth}
\addtolength{\gb}{-5.25pt}
\ifthenelse{\equal{\value{itemoddeven}}{0}}{%
\noindent\colorbox{bgeven}{\hskip-3pt\begin{minipage}{\gb}\Qitem[#1]{#2}\end{minipage}}%
\stepcounter{itemoddeven}%
}{%
\noindent\colorbox{bgodd}{\hskip-3pt\begin{minipage}{\gb}\Qitem[#1]{#2}\end{minipage}}%
\setcounter{itemoddeven}{0}%
}
}

\newcommand{\rawcode}[1]{\texttt{#1}}
\newcommand{\code}[2]{\mintinline{#1}{#2}}
\newcommand{\file}[1]{\texttt{\underline{#1}}}

\renewcommand{\theFancyVerbLine}{\tiny \texttt{\textcolor{gray}{\arabic{FancyVerbLine}}}}

\newcommand{\inputcode}[2]{\begin{mdframed}\inputminted[fontsize=\footnotesize,linenos,numbersep=4pt]{#1}{#2}\end{mdframed}}

\newenvironment{solution}{%
    \begin{mdframed}%
        \noindent\textbf{Solu\c{c}\~ao:}\par\medskip
}{%
    \end{mdframed}%
}

\newcommand{\inputsnippet}[4]{\begin{mdframed}[shadow=true,shadowsize=3pt,shadowcolor=gray]\inputminted[fontsize=\footnotesize,linenos,firstline=#2,lastline=#3,numbersep=2pt]{#1}{#4}\end{mdframed}}

\newcommand{\inputsyntax}[2]{\inputminted[fontsize=\footnotesize]{#1}{#2}}



\begin{document}

\begin{center}
	\includegraphics[scale=0.45]{fga}
	% fga.ps: 1087x103 pixel, 72dpi, 38.35x3.63 cm, bb=0 0 1087 103
\end{center}


\vspace{0.1in}

\begin{center}
    \textbf{\huge Paradigmas de Programação}

    \vspace{0.1in}
    \textbf{\Large Programação Funcional: Lista 1}
\end{center}\vskip1em

\renewcommand{\Qq}[1]{{#1}}

\section*{Parte A - Funções}

\Qitem{ \Qq{Usando apenas as funções uma ou mais dentre as funções \code{haskell}{head}, 
        \code{haskell}{tail}, \code{haskell}{take}, \code{haskell}{drop}, \code{haskell}{length} e
        \code{haskell}{null}, e os recursos e elementos presentes nas duas primeiras aulas sobre
        programação funcional, implemente as funções abaixo:

    \begin{Qlist}
        \item[(a)] \code{haskell}{second :: [a] -> Maybe a}, que recebe uma lista como parâmetro e 
            retorna o segundo elemento da lista, se existir. Ex.:
                \inputsyntax{haskell}{codes/second.hs}

        \begin{solution}
        O padrão \code{haskell}{(_ : x : _)} só casa com listas de pelo menos dois elementos e
        captura o segundo em \code{haskell}{x}; qualquer outra lista (vazia ou unitária) cai no
        segundo caso.

        \inputsnippet{haskell}{1}{4}{codes/sol/solutions.hs}
        \end{solution}

        \item[(b)] \code{haskell}{body :: [a] -> Maybe [a]}, que recebe uma lista como parâmetros e 
            retorna todos os elementos, exceto o primeiro e o último, se possível. Ex.:
                \inputsyntax{haskell}{codes/body.hs}

        \begin{solution}
        Listas vazias e unitárias não têm primeiro e último elementos distintos, então o
        resultado é \code{haskell}{Nothing}. Nos demais casos, o padrão \code{haskell}{(_:xs)}
        descarta o primeiro elemento, e \code{haskell}{take (length xs - 1) xs} descarta o último.
        Para uma lista de dois elementos, o resultado é \code{haskell}{Just []}.

        \inputsnippet{haskell}{6}{10}{codes/sol/solutions.hs}
        \end{solution}

        \item[(c)] \code{haskell}{median :: [a] -> Maybe a}, que retorna o elemento central da lista de
            $N$ elementos, se N é ímpar, ou o $N/2$-ésimo elemento, se $N$ é par, se possível. Ex.: 

                \inputsyntax{haskell}{codes/median.hs}

        \begin{solution}
        Em ambos os casos o elemento procurado é o de posição $\lfloor (N+1)/2\rfloor$
        (contando a partir de 1): se $N$ é ímpar, é o elemento central; se $N$ é par, é o
        $N/2$-ésimo. Como \code{haskell}{drop} descarta os $k$ primeiros elementos, descartamos
        $\lfloor (N+1)/2\rfloor - 1$ elementos e pegamos o \code{haskell}{head} do resto.

        \inputsnippet{haskell}{12}{15}{codes/sol/solutions.hs}
        \end{solution}
    \end{Qlist}
} }

\Qitem{ \Qq{Utilizando \textbf{obrigatoriamente} recursão, e opcionalmente guardas e os recursos 
        apresentados nas duas primeiras aulas sobre programação funcional, implemente as funções 
        abaixo:

    \begin{Qlist}
        \item[(a)] \code{haskell}{parity :: Int -> Int}, que recebe um inteiro positivo $n$ e retorna
            0 se $n$ é par, ou 1, se $n$ é ímpar.

        \begin{solution}
        Subtraindo $2$ de $n$ não se altera a paridade; a recursão termina em $0$ (par) ou
        $1$ (ímpar), que são exatamente os valores a retornar.

        \inputsnippet{haskell}{17}{21}{codes/sol/solutions.hs}
        \end{solution}

        \item[(b)] \code{haskell}{remainder :: Int -> Int -> Int}, que recebe dois inteiros positivos
            $a$ e $b$ e retorna o resto da divisão euclidiana de $a$ por $b$.

        \begin{solution}
        Se $a<b$, o resto é o próprio $a$; caso contrário, o resto de $a$ por $b$ é igual ao
        resto de $a-b$ por $b$.

        \inputsnippet{haskell}{23}{27}{codes/sol/solutions.hs}
        \end{solution}

        \item[(c)] \code{haskell}{quotient :: Int -> Int -> Int}, que recebe dois inteiros positivos
            $a$ e $b$ e retorna o quociente da divisão euclidiana de $a$ por $b$.

        \begin{solution}
        Se $a<b$, o quociente é $0$; caso contrário, é uma unidade a mais que o quociente de
        $a-b$ por $b$ (cabe mais um $b$ em $a$).

        \inputsnippet{haskell}{29}{33}{codes/sol/solutions.hs}
        \end{solution}

        \item[(d)] \code{haskell}{greater_common_divisor :: Int -> Int -> Int}, que recebe dois inteiros 
            positivos $a$ e $b$ e retorna o maior divisor comum $(a, b)$ de $a$ e $b$.
            Dica: $(a, b) = (b, a - b)$ e $(a, 0) = a$.

        \begin{solution}
        Caso base: $(a,0)=a$. Passo recursivo: o algoritmo de Euclides, $(a,b)=(b,a \bmod b)$,
        que é a dica $(a,b)=(b,a-b)$ aplicada repetidamente até que $a<b$. O resto é calculado
        pela função \code{haskell}{remainder} do item anterior.

        \inputsnippet{haskell}{35}{39}{codes/sol/solutions.hs}
        \end{solution}

        \item[(e)] \code{haskell}{semifactorial :: Int -> Int}, que recebe um inteiror positivo e
            retorna
        \[
            \mbox{\code{haskell}{semifactorial} $n$} = n(n - 2)(n - 4)\ldots 1
        \]
        e \code{haskell}{semifactorial} $1 = $ \code{haskell}{semifactorial} $0 = 1$.

        \begin{solution}
        Os casos base são $n=0$ e $n=1$ (o produto é $1$); para $n\geq 2$,
        $n!! = n\times (n-2)!!$.

        \inputsnippet{haskell}{41}{45}{codes/sol/solutions.hs}
        \end{solution}

        \item[(f)] \code{haskell}{power :: Int -> Int -> Int}, que recebe dois inteiros positivos
            $a$ e $n$ e retorna $a^n$.

        \begin{solution}
        $a^0=1$ e $a^n = a\times a^{n-1}$.

        \inputsnippet{haskell}{47}{51}{codes/sol/solutions.hs}
        \end{solution}

        \item[(g)] \code{haskell}{sum_of_first_n :: Int -> Int}, que recebe um natural $n$, e retorna 
            a soma dos primeiros $n$ números naturais.

        \begin{solution}
        A soma de zero termos é $0$, e $\sum_{i=1}^{n} i = n + \sum_{i=1}^{n-1} i$.

        \inputsnippet{haskell}{53}{57}{codes/sol/solutions.hs}
        \end{solution}

        \item[(h)] \code{haskell}{sum_of_first_n_squares :: Int -> Int}, que recebe um natural $n$, e 
            retorna a soma dos quadrados dos primeiros $n$ números naturais.

        \begin{solution}
        Análogo ao item anterior: $\sum_{i=1}^{n} i^2 = n^2 + \sum_{i=1}^{n-1} i^2$.

        \inputsnippet{haskell}{59}{63}{codes/sol/solutions.hs}
        \end{solution}

        \item[(i)] \code{haskell}{arithmetic_progression_sum :: Int -> Int -> Int -> Int}, que recebe 
            três naturais $a, r$ e $n$ e retorna a soma dos primeiros $n$ termos
            de uma progressão aritmética cujo primeiro termo é igual a $a$ e a razão é igual a $r$.

        \begin{solution}
        A soma dos $n$ termos é o primeiro termo $a$ mais a soma dos $n-1$ termos da
        progressão que começa em $a+r$ com a mesma razão.

        \inputsnippet{haskell}{65}{69}{codes/sol/solutions.hs}
        \end{solution}

        \item[(j)] \code{haskell}{geometric_progression_sum:: Int -> Int -> Int -> Int}, que recebe
            três naturais $a, q$ e $n$ e retorna a soma dos primeiros $n$ termos
            de uma progressão geométrica cujo primeiro termo é igual a $a$ e a razão é igual a $q$.

        \begin{solution}
        Como $a + aq + aq^2 + \cdots + aq^{n-1} = a + q\,(a + aq + \cdots + aq^{n-2})$, a soma de
        $n$ termos é $a$ mais $q$ vezes a soma de $n-1$ termos da mesma progressão.

        \inputsnippet{haskell}{71}{75}{codes/sol/solutions.hs}
        \end{solution}

        \item[(k)] \code{haskell}{binomial :: Int -> Int -> Int}, que recebe dois inteiros não-negativos
            $n$ e $m$, com $n \geq m$, e computa o coeficiente binomial $\binom{n}{m}$.

        \begin{solution}
        Usamos a relação de Pascal, $\binom{n}{m}=\binom{n-1}{m}+\binom{n-1}{m-1}$, com
        casos base $\binom{n}{n}=\binom{n}{0}=1$. A primeira guarda ($m>n$, que retorna $0$)
        não é alcançada para entradas válidas; serve apenas como proteção.

        \inputsnippet{haskell}{77}{82}{codes/sol/solutions.hs}
        \end{solution}

        \item[(l)] \code{haskell}{count_vowels :: String -> Int}, que recebe uma string de caracteres
            minúsculos $s$ e retorna o número de vogais (``$aeiou$'') presentes em $s$.

        \begin{solution}
        Uma string vazia tem $0$ vogais. Caso contrário, somamos $1$ se o primeiro caractere é
        uma vogal e continuamos a contagem no restante da string.

        \inputsnippet{haskell}{84}{89}{codes/sol/solutions.hs}
        \end{solution}
    \end{Qlist}
} }

\end{document}
